\documentclass[a4paper, serif, 10pt]{moderncv}

%% moderncv
% color and style
\moderncvcolor{black}
\moderncvstyle{banking}

%% page layout/margins
\usepackage[top=2.5cm, bottom=2.5cm, left=1.5cm, right=1.5cm]{geometry}

\nopagenumbers{}

%% Babel config for German language
% ref:
% - https://latex3.github.io/babel/guides/locale-german.html
\usepackage[ngerman]{babel}

%% fontspec
\usepackage{fontspec}

\IfFontExistsTF{Linux Libertine}{
  \setmainfont[Ligatures={TeX, Common}, Numbers=OldStyle]{Linux Libertine}
}{}
\IfFontExistsTF{Linux Libertine O}{
  \setmainfont[Ligatures={TeX, Common}, Numbers=OldStyle]{Linux Libertine O}
}{}

%% CTeX
% CJK support, used to show CJK characters in the resume
%
% - heading=false: disable ctex heading
% - punct=kaiming: use kaiming punctuations styles for CJK
% - scheme=plain: use plain scheme, do not override \normalsize font size
% - space=auto: space settings for CJK characters
%
% CJK fonts are resolved per locale language, e.g. Noto Serif/Sans CJK JP
% for Japanese resumes, with fallback to the next region if not installed
%
% ref:
% - http://ctan.mirrorcatalogs.com/language/chinese/ctex/ctex.pdf
\usepackage[UTF8, heading=false, punct=kaiming, scheme=plain, space=auto]{ctex}

\IfFontExistsTF{Noto Serif CJK SC}{
  \setCJKmainfont{Noto Serif CJK SC}
}{}
\IfFontExistsTF{Noto Sans CJK SC}{
  \setCJKsansfont{Noto Sans CJK SC}
}{}

%% line spacing
\usepackage{setspace}
\setstretch{1.125}

%% URL styling - use normal text instead of monospace
\urlstyle{same}

%% Patch \httplink and \httpslink to handle full URLs
% The original moderncv macros always prepend http:// or https:// to URLs.
% This causes issues when the URL already has a protocol (e.g., https://example.com)
% resulting in malformed URLs like https://https://example.com.
% This patch checks if the URL already contains "://" and uses it directly if so.
% Note: We use \str_set:Nx (x-expansion) to expand macros like \@homepage first.
\ExplSyntaxOn
\str_new:N \g_yamlresume_url_str

\RenewDocumentCommand{\httpslink}{O{}m}{
  \str_set:Nx \g_yamlresume_url_str {#2}
  \str_if_in:NnTF \g_yamlresume_url_str {://}
    { \url{#2} }
    { \url{https://#2} }
}

\RenewDocumentCommand{\httplink}{O{}m}{
  \str_set:Nx \g_yamlresume_url_str {#2}
  \str_if_in:NnTF \g_yamlresume_url_str {://}
    { \url{#2} }
    { \url{http://#2} }
}
\ExplSyntaxOff

\name{Lena Hartmann}{}
\title{Senior IT-Sicherheitsingenieurin}
\phone[mobile]{+49 30 12345678}
\email{lena.hartmann@security-mail.de}
\homepage{https://www.lhartmann.dev}

\address{Friedrichstraße 123, Berlin, Berlin, Deutschland, 10117}{}{}

\extrainfo{{\small \faGithub}\ \href{https://github.com/lhartmann}{@lhartmann} {} {} {} • {} {} {} 
{\small \faLinkedin}\ \href{https://www.linkedin.com/in/lhartmann}{@lhartmann}}

\begin{document}

\maketitle

\section{Basisinformationen}

\cvline{}{\begin{itemize}
\item Mehr als 7 Jahre Erfahrung in der IT-Sicherheit mit Fokus auf Angriffserkennung, Incident Response und Sicherheitsarchitekturen
\item Fundierte Kenntnisse in Penetrationstests, Schwachstellenmanagement und der Absicherung von Cloud-Infrastrukturen
\item Erfahrung mit SIEM-Systemen (Splunk, Wazuh), Firewalls und IDS/IPS sowie der Automatisierung von Sicherheitsprozessen
\item Teamorientiert, analytisch und kommunikationsstark; gewohnt, komplexe Sicherheitsrisiken verständlich zu erklären
\end{itemize}}
\section{Ausbildung}

\cventry{Okt. 2016–Sept. 2018}
        {Master, IT-Sicherheit, Punkte: 1.3}
        {Technische Universität Berlin}
        {\url{https://www.tu.berlin}}
        {}
        {\begin{itemize}
\item Schwerpunkt auf angewandter Kryptographie und der Absicherung verteilter Systeme
\item Masterarbeit zum Thema "Erkennung von Advanced Persistent Threats mittels Machine Learning"
\item Praktische Erfahrung durch Laborprojekte zu Penetrationstests und Sicherheitsarchitekturen
\end{itemize}
\textbf{Kurse}: Kryptographie und Sicherheitsprotokolle, Netzwerksicherheit und Firewalls, Secure Software Engineering, Incident Response und Forensik, Datenschutz und Compliance, Cloud Security und Virtualisierung}

\cventry{Okt. 2012–Sept. 2016}
        {Bachelor, Informatik, Punkte: 1.7}
        {Hochschule für Technik und Wirtschaft Berlin}
        {\url{https://www.htw-berlin.de}}
        {}
        {\begin{itemize}
\item Grundlagen in Programmierung, Datenbanken und Netzwerktechnik
\item Wahlpflichtmodule zu IT-Sicherheit und Systemadministration
\end{itemize}}
\section{Berufserfahrung}

\cventry{Apr. 2020–Heute}
        {Senior IT-Sicherheitsingenieurin}
        {SAP SE}
        {\url{https://www.sap.com}}
        {}
        {\begin{itemize}
\item Leitung von Penetrationstests und Sicherheitsaudits für cloudbasierte Unternehmensanwendungen
\item Entwicklung und Betrieb von SIEM-Regeln zur Erkennung von Angriffsmustern und Anomalien
\item Durchführung von Incident-Response-Maßnahmen und forensischen Analysen bei Sicherheitsvorfällen
\item Mitwirkung bei der Implementierung und Zertifizierung nach ISO 27001 und BSI-Grundschutz
\item Automatisierung von Sicherheitsprüfungen mit Python und Terraform
\end{itemize}
\textbf{Schlüsselwörter}: Penetration Testing, SIEM, Incident Response, ISO 27001, Cloud Security}

\cventry{Okt. 2018–März 2020}
        {IT-Sicherheitsanalystin}
        {Deutsche Telekom AG}
        {\url{https://www.telekom.com}}
        {}
        {\begin{itemize}
\item Überwachung und Analyse von Sicherheitsereignissen im 24/7-Security-Operations-Center
\item Bearbeitung von Sicherheitsvorfällen nach definierten Playbooks und Eskalationsprozessen
\item Durchführung von Schwachstellenanalysen und Bewertung von Risiken in Netzwerk- und Systemumgebungen
\item Erstellung von Sicherheitsreports und Präsentationen für technische und nicht-technische Stakeholder
\end{itemize}
\textbf{Schlüsselwörter}: SOC, SIEM, Schwachstellenanalyse, Netzwerksicherheit, Incident Response}
\section{Sprachen}

\cvline{Deutsch}{Muttersprache oder zweisprachig \hfill \textbf{Schlüsselwörter}: Muttersprache}
\cvline{Englisch}{Verhandlungssichere Kenntnisse \hfill \textbf{Schlüsselwörter}: Verhandlungssicher, TOEFL 110}
\cvline{Französisch}{Gute Kenntnisse \hfill \textbf{Schlüsselwörter}: Grundkenntnisse}
\section{Fähigkeiten}

\cvline{Netzwerksicherheit}{Experte \hfill \textbf{Schlüsselwörter}: Firewalls, VPN, IDS/IPS, Wireshark, Netzwerksegmentierung}
\cvline{Penetration Testing}{Erfahren \hfill \textbf{Schlüsselwörter}: Metasploit, Burp Suite, Nmap, OpenVAS, Nessus}
\cvline{Incident Response}{Erfahren \hfill \textbf{Schlüsselwörter}: SIEM (Splunk, Wazuh), Forensik, Threat Hunting, Playbook-Entwicklung}
\cvline{Cloud Security}{Fortgeschritten \hfill \textbf{Schlüsselwörter}: AWS, Azure, IAM, Kubernetes, Terraform}
\cvline{Compliance und Datenschutz}{Erfahren \hfill \textbf{Schlüsselwörter}: ISO 27001, BSI-Grundschutz, DSGVO, Risikomanagement}
\cvline{Programmierung und Automatisierung}{Erfahren \hfill \textbf{Schlüsselwörter}: Python, Bash, PowerShell, Git, CI/CD}
\section{Auszeichnungen}

\cventry{Nov. 2022}
        {Bundesverband IT-Sicherheit e.V.}
        {Cybersecurity Excellence Award}
        {}
        {}
        {\begin{itemize}
\item Auszeichnung für die erfolgreiche Einführung eines unternehmensweiten Threat-Intelligence-Programms
\item Anerkennung für besonderes Engagement in der Nachwuchsförderung im Bereich IT-Sicherheit
\end{itemize}}
\section{Zertifikate}

\cventry{Juni 2021}
        {ISC2}
        {Certified Information Systems Security Professional (CISSP)}
        {\url{https://www.isc2.org}}
        {}
        {}

\cventry{Nov. 2019}
        {Offensive Security}
        {Offensive Security Certified Professional (OSCP)}
        {\url{https://www.offsec.com}}
        {}
        {}

\cventry{März 2020}
        {TÜV Rheinland}
        {ISO/IEC 27001 Lead Auditor}
        {\url{https://www.tuv.com}}
        {}
        {}
\section{Veröffentlichungen}

\cventry{Dez. 2019}
        {Erkennung von Advanced Persistent Threats mittels Machine Learning}
        {Zeitschrift für IT-Sicherheit}
        {\url{https://www.zeitschrift-it-sicherheit.de/artikel/apt-ml}}
        {}
        {\begin{itemize}
\item Analyse von Netzwerk- und Hostdaten zur frühzeitigen Erkennung komplexer Angriffsmuster
\item Vergleich verschiedener Machine-Learning-Verfahren hinsichtlich Genauigkeit und False-Positive-Rate
\item Praxisnahe Evaluation in einer virtualisierten Unternehmensumgebung
\end{itemize}}
\section{Projekte}

\cventry{Jan. 2021–Dez. 2021}
        {Zentraler Aufbau eines SOC zur Erkennung und Reaktion auf Sicherheitsvorfälle in einer hybriden Cloud-Umgebung.}
        {Aufbau eines Security Operations Center (SOC)}
        {}
        {}
        {\begin{itemize}
\item Konzeption und Implementierung eines SOC für die zentrale Überwachung von Sicherheitsereignissen
\item Integration von SIEM-, Ticketing- und Threat-Intelligence-Systemen
\item Definition von Alarmierungs- und Eskalationsprozessen
\end{itemize}
\textbf{Schlüsselwörter}: SOC, SIEM, Threat Intelligence, Incident Response}

\cventry{März 2019–Feb. 2020}
        {Eine Plattform zur Sammlung, Normalisierung und Priorisierung von Threat-Intelligence-Daten aus verschiedenen Quellen.}
        {Threat-Intelligence-Plattform}
        {}
        {}
        {\begin{itemize}
\item Entwicklung einer Plattform zur Aggregation und Auswertung von Threat-Intelligence-Feeds
\item Automatisierte Anreicherung von Sicherheitsereignissen mit Kontextdaten
\item Bereitstellung von Dashboards für das Sicherheitsteam
\end{itemize}
\textbf{Schlüsselwörter}: Threat Intelligence, Python, APIs, Elasticsearch}
\section{Interessen}

\cvline{Capture the Flag (CTF)}{HackTheBox, TryHackMe, CTF-Wettbewerbe}
\cvline{Open Source}{Security-Tools, GitHub, Community}
\cvline{Sport}{Laufen, Klettern, Radfahren}
\section{Ehrenamtliche Tätigkeiten}

\cventry{Jan. 2021–Heute}
        {Mentorin für Nachwuchskräfte}
        {Bundesverband IT-Sicherheit e.V.}
        {\url{https://www.bv-it-sicherheit.de}}
        {}
        {\begin{itemize}
\item Betreuung von Studierenden und Berufseinsteigern bei CTF-Wettbewerben und Sicherheitsprojekten
\item Durchführung von Workshops zu Themen wie Penetration Testing und Incident Response
\item Förderung des Austauschs zwischen Hochschulen und Industrie
\end{itemize}}
    
\end{document}